In this appendix, we give formal proofs of the assertions connecting diffusion, denoising, and entropy encapsulated in \Cref{thm:diffusion-and-entropy} and \Cref{thm:conditioning_reduces_entropy}.
These results rely on the theory of \textit{Markov diffusion processes}, an extremely rich area of mathematics that blends probability, analysis, and geometry, and which has recently been exploited as the foundation for diffusion generative models, as we have mentioned.
For a general reference to this area of mathematics, consult \cite{Bakry2016-pl}.

\section{Diffusion Process}

\Cref{thm:diffusion-and-entropy} concerns the following setting: we observe a random variable $\vz_{\sigma} = \vz_o + \sigma \vg$, where $\vz_o$ is a random variable whose distribution $\mu$ has finite variance. Note that $\mu$ itself need not be a density: for example, $\vz_o$ could be a discrete random variable in everything that follows.

The following \textit{continuous time} perspective is helpful in analyzing the relationship between the entropy of $\vz_{\sigma}$ and $\vz$: consider an auxiliary time variable $t \in [0, \sigma]$, and a random variable $\vz_t$ defined analogously.
%It turns out (see e.g.\ \cite{Bakry2016-pl}) that the distribution of $\vz_{\sigma}$ is the same as the law of the (unique, strong) solution to the following \textit{stochastic differential equation} (SDE), evaluated at time $t =\sigma$:
%\begin{align*}
%    \mathrm{d} \vz_t &= \mathrm{d} \vw_t, \\ 
%    \vz_0 &\sim \vz_o.
%\end{align*}
%In this SDE, $\vw_t$ is a Brownian motion: roughly speaking, an uncountably-infinite collection of independent standard Gaussian random variables. 
By properties of Gaussian smoothing, for every $t > 0$, the random variable $\vz_t$ has a density $p_t$, which can be expressed analytically as $p_t = \varphi_{t} \mathop{\ast} \mu$, where $\varphi_t$ is the density of the standard Gaussian random variable with zero mean and covariance $t^2 \vI$.
A very fundamental fact about the density $p_t$ is that it satisfies the \textit{heat equation}:\footnote{See \cite{Bakry2016-pl}, and note that this correspondence is at the heart of modern diffusion generative models, although we omit these probabilistic connections here for simplicity.} specifically, the PDE
\begin{equation*}
    \partial_t p_t = \frac{1}{2} \Delta p_t, 
\end{equation*}
where $\partial_t$ denotes differnetiation in time, and $\Delta = \sum_{i=1}^d \partial^2_{\vx_i}$ denotes the Laplacian operator.
By exploiting the heat equation, we find a very simple proof that \textit{entropy increases as $t$ increases}:
calculate
\begin{equation}\label{eq:de-bruijn}
  \begin{split}
  \partial_t h(\vz_t)
  &=
  -\frac{1}{2} \int (1 + \log p_t) \Delta p_t
  \\
  &=
  \frac{1}{2} \int \ip*{\nabla \log p_t}{\nabla p_t}
  \\
  &= \frac{1}{2} \int \frac{\norm{\nabla p_t}_2^2}{p_t}
  \geq 0,
  \end{split}
\end{equation}
where the first line uses the chain rule, the heat equation, and the expression $h(\vz_t) = \int p_t \log p_t$; and the second line uses an integration by parts.
From this calculation, \Cref{thm:diffusion-and-entropy} immediately follows.

\section{Denoising Process}

\Cref{thm:conditioning_reduces_entropy} concerns the corresponding ``inverse''
process: we start with the noisy random variable $\vz_{\sigma}$, and we seek to
compare its entropy to that of the \textit{denoised} random variable
\begin{equation*}
  \vmu_{\sigma}(\vz_{\sigma}) = \bE[ \vz_o \mid \vz_{\sigma} ] = \vz_{\sigma}
  + \sigma^2 \nabla \log p_{\sigma}(\vz_{\sigma}),
\end{equation*}
where we applied Tweedie's formula (\Cref{thm:tweedie}).
This can be seen as a perturbation of the distribution of the random variable
$\vz_{\sigma}$ by the \textit{score function vector field}. 
% We can
% characterize how this perturbation affects the entropy of the distribution
% by exploiting a deep connection between stochastic differential equations and
% diffusion-denoising processes, 
How does this perturbation affect the entropy of the noisy observation? To
answer this question, we develop a continuous time 


